\begin{abstract}
LLM-based multi-agent systems compress conversation history into memory summaries to manage long-horizon context.
We show this compression step opens a previously unexamined safety blind spot: toxic influence can be \emph{laundered} through summarization, producing memory states that classifiers mark as clean but that still steer downstream agent behavior toward hostility.
% Standard per-turn toxicity filtering, the dominant safety paradigm for multi-agent systems, structurally cannot detect this contamination because no individual artifact crosses the detection threshold.
Therefore, we first identify and characterize \emph{memory laundering} as a distinct safety failure mode: across $200$ paired seeds, $100\%$ of toxic-condition memory summaries score below the standard Detoxify threshold $\tau{=}0.5$ (and below $\tau{=}0.03$), yet conditioning on them shifts downstream Detoxify toxicity (in $[0,1]$ units) by $0.139$ relative to a matched neutral baseline ($p{=}3.75{\times}10^{-8}$).
% The impact: any safety monitor that filters per-turn classifier scores will deem these systems safe while behavioral contamination accumulates.
We then introduce the \emph{sub-threshold propagation gap} (SPG), the first metric to quantify behavioral influence in classifier-clean states by measuring the downstream toxicity difference conditioned on $\mathrm{tox}(M_t){<}\tau$.
Unlike mean toxicity, which averages across all states, SPG isolates exactly the regime a deployed monitor would label safe; it remains significantly positive at every threshold tested (all $p{<}0.001$), surfacing the laundering signal that mean-based metrics dilute.
Third, we propose \emph{state-aware mitigation}: read-side context sanitization and write-side state gating.
Through paired interventions across seven graph topologies ($1{,}400$ rollouts), we disentangle the transcript and memory channels where transcript is the dominant overt channel ($\Delta\mu{=}0.13$ single-injection, $0.38$ multi-injection in Detoxify units; $p{<}0.01$), and memory carries the hidden sub-threshold channel. And show that sanitization applied \emph{before} compression eliminates SPG ($\mathrm{SPG}{\approx}0.0004$) while sanitization applied \emph{after} compression leaves it intact ($\mathrm{SPG}{\approx}0.086$).
% The impact: laundering is an order-of-operations property, and any defense must intercept toxic content before it enters the summarizer, not after.
Together, these results reframe agent safety as a state-control problem and identify a specific architectural constraint that per-turn filtering and parameter-level unlearning (DPO) alone cannot satisfy.
\end{abstract}
% \begin{abstract}
% When LLM agents compress conversation history into memory summaries, toxic influence can be \emph{laundered}: the summary's toxicity score drops over $40\times$ (from $0.910$ to $0.022$), falling well below any standard classifier threshold, yet downstream agents conditioned on this memory still produce elevated toxicity across three subsequent turns.
% We formalize this as \emph{memory laundering} and show it constitutes a blind spot for per-turn safety evaluation: in $62\%$ of toxic-condition memory states, hostile framing persists at classifier scores below $0.025$ while continuing to shift agent behavior.
% To detect this hidden influence, we introduce the \emph{sub-threshold propagation gap} (SPG)---the downstream toxicity difference conditioned only on memory states that classifiers mark as clean---and show it remains significantly positive even when no individual message exceeds the detection threshold.
% Through paired intervention experiments across seven graph topologies (248 seeds, 1{,}400+ rollouts), we disentangle two propagation channels: the raw transcript and compressed memory.
% The transcript is the dominant channel ($\Delta\mu \approx 0.13$ single-injection, $0.38$ multi-injection; $p{<}0.01$), but memory creates a hidden secondary channel that evades standard detection.
% %TODO: add one sentence on channel isolation result once smoking-gun experiment is complete
% We show that neither per-turn output filtering nor parameter-level intervention (DPO) suffices alone, and propose \emph{state-aware mitigation}---read-side context sanitization and write-side state gating---as the necessary complement.
% %TODO: add combined-intervention result sentence once joint condition is run
% \end{abstract}



% \begin{abstract}
% When LLM agents compress conversation history into memory summaries, toxic influence can be laundered: the summary's toxicity score drops by over an order of magnitude, falling well below any standard classifier threshold, yet downstream agents conditioned on this memory still produce overtly hostile output. We formalize this as \emph{memory laundering} and show it constitutes a blind spot for per-turn safety evaluation. Across $200$ seeds, $100\%$ of toxic-condition memory summaries score below the standard threshold $\tau = 0.5$ --- and indeed below $\tau = 0.03$ --- while conditioning on them still shifts downstream toxicity by $0.139$ relative to the matched neutral baseline ($p = 3.75 \times 10^{-8}$). To detect this hidden influence, we introduce the \emph{sub-threshold propagation gap} (SPG) --- the downstream toxicity difference conditioned only on memory states classifiers mark as clean --- and show it remains significantly positive at every threshold tested (all $p < 0.001$). Through paired interventions across seven graph topologies ($200$ seeds, $1{,}400$ rollouts), we disentangle two propagation channels: the raw transcript and compressed memory. The transcript dominates overt propagation ($\Delta\mu = 0.27$ on chain), while memory carries hidden behavioral influence below classifier thresholds. The two channels exhibit distinct metric signatures: transcript shows up as $\Delta\mu$, memory shows up as SPG. We further show that threshold-based sanitization eliminates SPG when applied to source messages before compression ($\mathrm{SPG} \approx 0.0004$) but leaves it essentially intact when applied to summaries after compression ($\mathrm{SPG} \approx 0.086$), establishing laundering as a property of the order of operations relative to summarization rather than of threshold sanitization itself. The phenomenon generalizes beyond chain to branching topologies (tree $\mathrm{SPG} = 0.049$, $p < 0.001$). Our findings reframe agent safety as a state-control problem and identify an architectural constraint any defense must respect: sanitization must occur before, not after, the compression step.
% \end{abstract}